% ============================================================================
% Մանկավարժական լուծումներ - Դիսկրետ մաթեմատիկա, Գործնական 7
% (Բազմությունների տրոհումներ և կցում-արտաքսում). Խնդիրներ 1, 7, 10, 12.
% Հայերեն տարբերակ, XeLaTeX (fontspec + Sylfaen). Կազմել՝ xelatex-ով երկու անգամ։
% Տերմինաբանությունը՝ glossary_discrete_focused.md / glossary_Apr.csv-ի համաձայն։
% ============================================================================
\documentclass[11pt,a4paper]{article}

% ---- Տառատեսակ (XeLaTeX) ----
\usepackage{fontspec}
\setmainfont{Sylfaen}

% Հայերենը plain XeLaTeX-ում տողադարձման կանոններ չունի, ուստի երկար տողերը չեն
% կարող կոտրվել և դուրս կգան լուսանցքից։ Թող TeX-ը ձգի բառամիջյան բացատները։
\tolerance=2000
\emergencystretch=3.5em

% ---- Մաթեմատիկա ----
\usepackage{amsmath,amssymb,amsthm}
\usepackage{mathtools}

% ---- Դասավորություն ----
\usepackage[margin=2.3cm]{geometry}
\usepackage{enumitem}
\usepackage{booktabs}
\usepackage{array}
\usepackage{colortbl}

% ---- Գրաֆիկա ----
\usepackage{xcolor}
\usepackage[most]{tcolorbox}
\usepackage{tikz}
\usetikzlibrary{arrows.meta,positioning,calc,fit,backgrounds,shapes.misc}
\usepackage{pgfplots}
\pgfplotsset{compat=1.18}
\usepackage{hyperref}
\hypersetup{colorlinks=true,linkcolor=blue!55!black,urlcolor=blue!55!black}

% ---- Գունապնակ (Հայաստանի դրոշի գույները) ----
\definecolor{armRed}{HTML}{D90012}
\definecolor{armBlue}{HTML}{0033A0}
\definecolor{armOrange}{HTML}{F2A800}
\definecolor{ink}{HTML}{22303C}

\setlength{\parindent}{0pt}
\setlength{\parskip}{5pt}

% ---- Վանդակներ ----
\newtcolorbox{problembox}[1][]{enhanced, breakable, sharp corners=south,
  colback=armOrange!10, colframe=armOrange!75!black, boxrule=0.8pt, arc=2mm,
  fonttitle=\bfseries\color{white}, coltitle=white,
  attach boxed title to top left={xshift=4mm,yshift=-2.4mm},
  boxed title style={colback=armOrange!85!black, arc=1mm},
  title={Խնդիր #1}}

\newtcolorbox{ideabox}{enhanced, breakable,
  colback=armBlue!6, colframe=armBlue!70!black, boxrule=0.8pt, arc=2mm,
  fonttitle=\bfseries\color{white}, coltitle=white,
  attach boxed title to top left={xshift=4mm,yshift=-2.4mm},
  boxed title style={colback=armBlue!75!black, arc=1mm},
  title={Գաղափարը}}

\newtcolorbox{methodbox}{enhanced, breakable,
  colback=gray!8, colframe=gray!55!black, boxrule=0.7pt, arc=2mm,
  fonttitle=\bfseries, title={$\bigstar$\ Հիշարժան մեթոդ}}

\newtcolorbox{answerbox}{enhanced, sharp corners=south,
  colback=green!8, colframe=green!50!black, boxrule=1pt, arc=2mm, left=4mm,right=4mm}
\newcommand{\answer}[1]{\begin{answerbox}\centering\textbf{Պատասխան՝}\quad #1\end{answerbox}}

% ---- Վերնագիր ----
\title{\vspace{-1.2cm}\textbf{Դիսկրետ մաթեմատիկա - Գործնական 7}\\[2pt]
\large Բազմությունների տրոհումներ և կցում-արտաքսում՝ մանրամասն լուծումներ\\[2pt]
\normalsize\itshape դատողությունները՝ քայլ առ քայլ բացված}
\author{}
\date{}

\begin{document}
\maketitle
\vspace{-1.0cm}

\section*{Նախապատրաստում}
Այս թերթիկում աշխատում են երկու գործիք։ Առաջինը \textbf{երկրորդ սեռի Ստիռլինգի թիվն} է $S(n,k)$՝
$n$-տարրանոց բազմությունը $k$ ոչ դատարկ չպիտակավորված դասերի տրոհելու տարբերակների քանակը։ Երկրորդը
\textbf{կցում-արտաքսման սկզբունքն (ԿԱՍ)} է՝ «վատ» հատկությունների ցանկից խուսափող տարրերը հաշվելու
համար սկսիր ընդհանուրից, հանիր յուրաքանչյուր վատ բազմություն, ապա հետ ավելացրու յուրաքանչյուր զույգ
հատում (որը երկու անգամ հանվեց), հանիր եռյակ հատումները և այդպես շարունակ։ Այս թերթիկի մեծ մասը
քողարկված ԿԱՍ է՝ բաժանելիության սահմանափակումներով թվերի հաշվումը, կամ վերին սահմաններով հավասարման
ամբողջ լուծումների հաշվումը, նույն նշանափոխ-գումարի գաղափարն է։ Լուծում ենք 1, 7, 10 և 12 խնդիրները։

% ============================================================================
\section*{Խնդիր 1 - Բազմության տրոհումը ոչ դատարկ դասերի}

\begin{problembox}[1]
Քանի՞ տարբերակով կարելի է $5$-տարրանոց բազմությունը տրոհել $3$ ոչ դատարկ դասերի։ (Դասերը
չպիտակավորված են՝ նշանակություն ունի միայն խմբավորումը, ոչ թե դասերի անունները կամ կարգը։)
\end{problembox}

\begin{ideabox}
Այս քանակը հենց երկրորդ սեռի Ստիռլինգի $S(5,3)$ թիվն է, և այն ճշտելու ամենամաքուր ձևը \textbf{մեկ
տարրի վրա կենտրոնանալով կառուցված ռեկուրենտ առնչությունն} է։ Առանձնացրու մեկ ֆիքսված տարր, ասենք
վերջինը՝ $x_5$։ $3$ ոչ դատարկ դասերի ցանկացած տրոհման մեջ $x_5$-ի հետ տեղի է ունենում ճիշտ երկու
բաներից մեկը՝
\begin{itemize}[nosep,leftmargin=2em]
  \item \textbf{$x_5$-ը մենակ է} իր սեփական դասում։ Այդ դեպքում մյուս $4$ տարրերը պետք է կազմեն
        մնացած $2$ ոչ դատարկ դասերը՝ $S(4,2)$ տարբերակ։
  \item \textbf{$x_5$-ը միանում է մյուսներին։} Այդ դեպքում մյուս $4$ տարրերն արդեն կազմում են բոլոր
        $3$ ոչ դատարկ դասերը ($S(4,3)$ տարբերակ), և $x_5$-ը գցվում է այդ $3$ դասերից մեկի մեջ՝
        $3\cdot S(4,3)$ տարբերակ։
\end{itemize}
Այս երկու դեպքերը փոխբացառող են և ընդգրկում են ամեն ինչ, ուստի դրանք գումարում ենք։ Սա
$S(n,k)=S(n-1,k-1)+k\,S(n-1,k)$ ընդհանուր ռեկուրենտ առնչությունն է, և այն թույլ է տալիս բարձրանալ
փոքր հայտնի արժեքներից։
\end{ideabox}

\textbf{Քայլ 1 - Ստանանք ավելի փոքր $S(4,2)$ և $S(4,3)$ արժեքները։}
Նույն ռեկուրենտ առնչությամբ (և $S(n,1)=1$, $S(n,n)=1$)՝
$$S(4,2)=S(3,1)+2\,S(3,2)=1+2\cdot 3=7,\qquad S(4,3)=S(3,2)+3\,S(3,3)=3+3\cdot1=6.$$
(Այստեղ $S(3,2)=S(2,1)+2S(2,2)=1+2=3$։)

\textbf{Քայլ 2 - Միավորենք։}
$$S(5,3)=S(4,2)+3\,S(4,3)=7+3\cdot 6=7+18=25.$$

\textbf{Ստուգում։} Անկախ բանաձևը
$S(n,k)=\frac{1}{k!}\sum_{j=0}^{k}(-1)^{j}\binom{k}{j}(k-j)^{n}$ տալիս է
$$S(5,3)=\frac{1}{6}\big(3^5-3\cdot 2^5+3\cdot 1^5\big)=\frac{243-96+3}{6}=\frac{150}{6}=25.$$
$\{1,2,3,4,5\}$-ի բոլոր տրոհումների ուղիղ համակարգչային թվարկումը գտնում է ճիշտ $25$-ը՝ երեք դասով։
\checkmark

\answer{$S(5,3)=25$}

\begin{methodbox}
\textbf{Փոխանցելի տեխնիկա - Ստիռլինգի ռեկուրենտ առնչությունը՝}
$$S(n,k)=\underbrace{S(n-1,k-1)}_{\text{նոր տարրը մենակ է}}+\underbrace{k\,S(n-1,k)}_{\text{նոր տարրը միանում է $k$ դասերից մեկին}}.$$
Մեկ տարր առանձնացնելը և «մենա՞կ է, թե՞ միանում է» հարցնելը ռեկուրսիվ կառուցվածքների հաշվման
հիմնական հնարքն է։ \textbf{Զգուշացում՝} եթե դասերը \emph{պիտակավորված} լինեին (տարբերակելի արկղեր),
քանակը կլիներ $k!\,S(n,k)$, իսկ դատարկ արկղեր թույլ տալու դեպքում՝ $k^{n}$։
\end{methodbox}

\medskip
\textbf{Աղյուսակի կառուցումը։} Յուրաքանչյուր վանդակ «վերին-ձախ հարևանն» է գումարած
$k\times$«ձախ հարևանը»՝ մեր թիրախ $S(5,3)=25$-ը գունանշված է։

\begin{center}
\small
\renewcommand{\arraystretch}{1.15}
\begin{tabular}{@{}c|cccccc@{}}
\toprule
$n \backslash\, k$ & 1 & 2 & 3 & 4 & 5\\
\midrule
1 & 1 &    &    &   & \\
2 & 1 & 1  &    &   & \\
3 & 1 & 3  & 1  &   & \\
4 & 1 & 7  & 6  & 1 & \\
5 & 1 & 15 & \cellcolor{armOrange!30}\textbf{25} & 10 & 1\\
\bottomrule
\end{tabular}
\end{center}

% ============================================================================
\section*{Խնդիր 7 - 11-ի վրա բաժանվող, բայց ոչ 5-ի և ոչ 3-ի վրա}

\begin{problembox}[7]
Հաշվիր $1$-ից $16500$ ընկած ամբողջ թվերը, որոնք բաժանվում են $11$-ի վրա, բայց \textbf{չեն} բաժանվում
$5$-ի վրա և \textbf{չեն} բաժանվում $3$-ի վրա։
\end{problembox}

\begin{ideabox}
«$11$-ի վրա բաժանվող» արտահայտությունը ֆիքսում է մեր տիեզերքը՝ մոռացիր մնացած բոլոր թվերը և նայիր
միայն տիրույթում գտնվող $11$-ի $1500$ բազմապատիկներին։ Այդ տիեզերքի ներսում պետք է դուրս նետենք
նրանց, որոնք \emph{նաև} բաժանվում են $5$-ի կամ $3$-ի վրա։ «Կամ»-ը \textbf{կցում-արտաքսման} ազդանշանն
է՝ եթե պարզապես գումարենք ($5$-ի բազմապատիկներ) և ($3$-ի բազմապատիկներ), ապա \emph{և՛} $5$-ի,
\emph{և՛} $3$-ի վրա բաժանվող թվերը կհաշվվեն երկու անգամ, ուստի դրանք մեկ անգամ հետ ենք ավելացնում։

Մաքուր հնարք, որը թվաբանությունը պահում է ճշգրիտ՝ «$11$-ի բազմապատիկ և $5$-ի բազմապատիկ» նշանակում է
«$55$-ի բազմապատիկ» (քանի որ $5,11$-ը փոխադարձաբար պարզ են), նմանապես $3$-ը տալիս է $33$, իսկ $5$-ը և
$3$-ը միասին՝ $165$։ Քանի որ $16500=11\cdot1500$ և $16500$-ը նաև $55$-ի, $33$-ի և $165$-ի բազմապատիկ
է, ստորև յուրաքանչյուր ստորին ամբողջ մաս ճշգրիտ բաժանում է՝ կլորացման մասին անհանգստանալու կարիք չկա։
\end{ideabox}

\textbf{Քայլ 1 - Տիեզերքը՝ $11$-ի բազմապատիկները։}
$$\left\lfloor\frac{16500}{11}\right\rfloor=1500.$$

\textbf{Քայլ 2 - «Վատ»-երը (նաև բաժանվում են $5$-ի կամ $3$-ի վրա)։}
$11$-ի բազմապատիկների մեջ՝
$$\#(\text{նաև }5)=\frac{16500}{55}=300,\quad
  \#(\text{նաև }3)=\frac{16500}{33}=500,\quad
  \#(\text{նաև }5\text{ և }3)=\frac{16500}{165}=100.$$

\textbf{Քայլ 3 - Կցում-արտաքսում «կամ»-ի համար։}
$11$-ի բազմապատիկներ, որոնք բաժանվում են $5$-ի \textbf{կամ} $3$-ի վրա՝
$$300+500-100=700.$$

\textbf{Քայլ 4 - Հանենք տիեզերքից։}
$$1500-700=800.$$

\textbf{Ստուգում։} $1,\dots,16500$ վրայով ուղիղ ցիկլը, որը հաշվում է $11\mid n,\ 5\nmid n,\ 3\nmid n$
պայմանին բավարարող $n$-երը, վերադարձնում է ճիշտ $800$։ Խաչաձև ստուգում կոտորակներով՝ $11$-ի
բազմապատիկների մեջ «ոչ $5$, ոչ $3$» պայմանին գոյատևող մասը $\frac{4}{5}\cdot\frac{2}{3}=\frac{8}{15}$
է, և $1500\cdot\frac{8}{15}=800$ (այս մաքուր կոտորակն աշխատում է, քանի որ $1500$-ը $15$-ի բազմապատիկ
է)։ \checkmark

\answer{$1500-(300+500-100)=800$}

\begin{methodbox}
\textbf{Փոխանցելի տեխնիկա - հաշվիր ճիշտ տիեզերքում, ապա ԿԱՍ՝} «$A$, բայց ոչ $B$ կամ $C$»-ի համար՝
սահմանափակվիր $A$-ով, ապա հանիր $|B\cup C|=|B|+|C|-|B\cap C|$ (ամենը չափված $A$-ի ներսում)։
Բաժանելիության համար «$a$-ի և $b$-ի բազմապատիկ» նշանակում է «$\operatorname{lcm}(a,b)$-ի
բազմապատիկ», և $N$-ից ոչ մեծ քանակը $\lfloor N/\operatorname{lcm}\rfloor$ է։
\end{methodbox}

\medskip
\textbf{Վեննի պատկեր ($11$-ի $1500$ բազմապատիկների ներսում քանակները)։} Ջնջում ենք գունանշված
համընկնող շրջանը՝ $300+500-100=700$, և պահում $800$-ը դրսում։

\begin{center}
\begin{tikzpicture}[font=\small]
  \draw[rounded corners=6pt, thick, armBlue!70!black] (-3.6,-2.1) rectangle (3.6,2.3);
  \node[armBlue!70!black] at (-2.3,1.95) {$11$-ի բազմ.՝ $1500$};
  \begin{scope}
    \clip (-0.9,0) circle (1.5);
    \fill[armRed!18] (0.9,0) circle (1.5);
  \end{scope}
  \draw[armRed!75!black, thick, fill=armRed!10, fill opacity=0.45] (-0.9,0) circle (1.5);
  \draw[armOrange!80!black, thick, fill=armOrange!18, fill opacity=0.45] (0.9,0) circle (1.5);
  \node at (-1.65,0) {$5$-ի վրա};
  \node at ( 1.65,0) {$3$-ի վրա};
  \node at (-1.45,-0.95) {$200$};
  \node at ( 1.45,-0.95) {$400$};
  \node at ( 0,-0.0) {$100$};
  \node at ( 0,-1.7) {պահված (ոչ մեկը)՝ $\mathbf{800}$};
\end{tikzpicture}
\end{center}

% ============================================================================
\section*{Խնդիր 10 - Եռանկյունները 8 մասի բաժանված եռանկյան մեջ}

\begin{problembox}[10]
$ABC$ եռանկյան յուրաքանչյուր կողմ բաժանված է $8$ հավասար մասի, և բաժանման կետերով տարված են բոլոր
ուղիղները՝ զուգահեռ երեք կողմերին, ստեղծելով եռանկյունային ցանց։ Հաշվիր բոլոր եռանկյունները, որոնց
գագաթները ցանցի հանգույցներ են ($A,B,C$ անկյունները նույնպես հաշվվում են որպես ցանցի հանգույցներ) և
որոնց երեք կողմերից յուրաքանչյուրը զուգահեռ է $ABC$-ի որևէ կողմին։ Սա ներառում է $ABC$-ի հետ նույն
ուղղվածությամբ եռանկյունները (\textbf{ուղիղ}) և շրջվածները (\textbf{շրջված})։
\end{problembox}

\begin{ideabox}
«$ABC$-ին զուգահեռ կողմեր» նշանակում է, որ մեր հաշված յուրաքանչյուր եռանկյուն նման է $ABC$-ին և ուղիղ
նստած է ցանցի վրա։ Հնարավոր են ընդամենը երկու ուղղվածություն՝ \textbf{ուղիղ} ($ABC$-ի պես վեր
ուղղված) և \textbf{շրջված} (շրջված/վար ուղղված)։ Ցանցն ունի փոքր ուղիղ միավոր եռանկյուններ և փոքր
շրջված միավոր եռանկյուններ, և ցանկացած ուղղվածության ավելի մեծ եռանկյունը պարզապես փոքրերի միավորում
է։ Ուստի պլանն է՝ հաշվել ամեն չափի ուղիղ եռանկյունները, հաշվել ամեն չափի շրջված եռանկյունները և
գումարել։

Ինչու՞ բաժանել ըստ ուղղվածության։ Որովհետև երկու ընտանիքները ենթարկվում են \emph{տարբեր} չափային
սահմանափակումների։ $k$ կողմ ունեցող ուղիղ եռանկյունը տեղավորվում է այնքան ժամանակ, քանի դեռ $k\le n$։
Բայց $k$ կողմ ունեցող շրջված եռանկյանը տեղ է պետք բոլոր կողմերից՝ այն տեղավորվում է միայն երբ $k$-ն
ամենաշատը $\lfloor n/2\rfloor$ է, ինչը պատճառն է, որ մեծ շրջված եռանկյունները պարզապես գոյություն
չունեն։ Երկու հաշիվներն առանձին պահելն է, ինչը խնդիրը կառավարելի է դարձնում։ Այստեղ $n=8$։
\end{ideabox}

\textbf{Քայլ 1 - $k$ կողմ ունեցող ուղիղ եռանկյունները։}
$k$ կողմ ունեցող ուղիղ եռանկյունը որոշվում է իր վերին գագաթի դիրքով։ Այդ գագաթը կարող է նստել ցանկացած
տեղ, որը դեռ տեղ է թողնում իր տակ գտնվող $k$-կողմ եռանկյան համար, և այդ դիրքերը կազմում են $n-k$ կողմ
ունեցող եռանկյունային զանգված, որը պարունակում է $1+2+\dots+(n-k+1)=\frac{(n-k+1)(n-k+2)}{2}$ կետ։
Յուրաքանչյուր չափի քանակն ինքնին բինոմիալ գործակից է՝ $\frac{(n-k+1)(n-k+2)}{2}=\binom{n-k+2}{2}$,
ուստի բոլոր չափերի վրայով գումարելը՝ $k=1,\dots,8$, տալիս է $\binom{2}{2}+\binom{3}{2}+\dots+\binom{9}{2}$։
«Հոկեյի փայտիկի» նույնությամբ $\bigl(\sum_{m=2}^{n+1}\binom{m}{2}=\binom{n+2}{3}\bigr)$ սա փլուզվում է
մեկ բինոմիալ գործակցի՝
$$\sum_{k=1}^{8}\frac{(n-k+1)(n-k+2)}{2}=\binom{n+2}{3}=\binom{10}{3}=120.$$
Անդամ առ անդամ՝ $36,28,21,15,10,6,3,1$ $k=1,\dots,8$-ի համար, և $36+28+21+15+10+6+3+1=120$։

\textbf{Քայլ 2 - $k$ կողմ ունեցող շրջված (վար ուղղված) եռանկյունները։}
$k$ կողմ ունեցող շրջված եռանկյուն գոյություն ունի միայն $k\le\lfloor 8/2\rfloor=4$-ի համար, և կա
$\frac{(n-2k+1)(n-2k+2)}{2}$ հատ $k$-կողմ՝
$$k=1:\ \tfrac{7\cdot 8}{2}=28,\quad k=2:\ \tfrac{5\cdot 6}{2}=15,\quad
  k=3:\ \tfrac{3\cdot 4}{2}=6,\quad k=4:\ \tfrac{1\cdot 2}{2}=1.$$
Շրջված եռանկյունների գումարը՝ $28+15+6+1=50$։

\textbf{Քայլ 3 - Ընդամենը։}
$$\underbrace{120}_{\text{ուղիղ}}+\underbrace{50}_{\text{շրջված}}=170.$$

\textbf{Ստուգում։} Եռանկյունային ցանցի վրայով ուղիղ համակարգչային թվարկումը (ստուգիր ցանցի կետերի
յուրաքանչյուր եռյակ, որը կազմում է ցանկացած չափի ուղիղ կամ շրջված եռանկյուն) վերադարձնում է $120$
ուղիղ և $50$ շրջված, ընդամենը՝ $170$։ $n$ կողմ ունեցող եռանկյան հայտնի փակ բանաձևը նույնպես
համընկնում է՝ $n=8$-ի համար $\big\lfloor \frac{n(n+2)(2n+1)}{8}\big\rfloor=\frac{8\cdot 10\cdot 17}{8}=170$։
\checkmark

\answer{$120\ (\text{ուղիղ})+50\ (\text{շրջված})=170$}

\medskip
\textbf{Կարևոր ընթերցման նշում։} Եթե նախատեսված հաշիվը լիներ միայն $ABC$-ի հետ \emph{նույն
ուղղվածությամբ} (ուղիղ) եռանկյունները, պատասխանը կլիներ $120$։ Պայմանը պահանջում է բոլոր ցանցային
եռանկյունները՝ $ABC$-ին զուգահեռ կողմերով, ինչը ներառում է շրջվածները, ուստի լրիվ ընդամենը՝
$\mathbf{170}$։ Ներկայացնում ենք բաժանումն այնպես, որ երկու ընթերցումներն էլ ծածկված լինեն։

\begin{methodbox}
\textbf{Փոխանցելի տեխնիկա - հաշվիր ըստ չափի և ըստ ուղղվածության՝} $n$ կողմ ունեցող եռանկյունային
ցանցում՝ ուղիղ եռանկյունների ընդհանուր քանակը $\binom{n+2}{3}$ է ($\frac{(n-k+1)(n-k+2)}{2}$-ի
գումարը)։ $k$ կողմ ունեցող շրջված եռանկյունները գոյություն ունեն միայն $2k\le n$-ի դեպքում և կազմում
են $\frac{(n-2k+1)(n-2k+2)}{2}$։ Երբ պատկերը կարող է հանդես գալ մեկից ավելի ուղղվածությամբ, նախ
բաժանիր հաշիվն ըստ ուղղվածության՝ ուղղվածությունները սովորաբար ունեն տեղավորման տարբեր
սահմանափակումներ։
\end{methodbox}

\medskip
\textbf{Ցանցը (այստեղ պարզության համար գծված $4$ կողմով)։} Ուղիղ եռանկյունները $ABC$-ի պես ուղղված են
վեր։ Գունանշված միակ եռանկյունը \emph{շրջված} է։ $k$ կողմ ունեցող շրջված եռանկյունը տեղավորվում է
միայն երբ $2k\le n$ (այն պետք է նստի $2k$ կողմ ունեցող ուղիղ եռանկյան ներսում), ինչը պատճառն է, որ
մեծ շրջվածներին տեղ չի մնում։

\begin{center}
\begin{tikzpicture}[scale=1.05, font=\small]
  \def\nn{4}
  \pgfmathsetmacro{\hh}{sqrt(3)/2}
  \foreach \j in {0,...,4}{
    \pgfmathsetmacro{\xs}{\j*0.5}
    \pgfmathsetmacro{\ys}{\j*\hh}
    \pgfmathsetmacro{\xe}{(\nn-\j)+\j*0.5}
    \draw[gray!55] (\xs,\ys) -- (\xe,\ys);
  }
  \foreach \i in {0,...,4}{
    \pgfmathsetmacro{\xs}{\i}
    \pgfmathsetmacro{\xe}{\i+(\nn-\i)*0.5}
    \pgfmathsetmacro{\ye}{(\nn-\i)*\hh}
    \draw[gray!55] (\xs,0) -- (\xe,\ye);
  }
  \foreach \i in {0,...,4}{
    \pgfmathsetmacro{\xs}{\i}
    \pgfmathsetmacro{\xe}{\i*0.5}
    \pgfmathsetmacro{\ye}{\i*\hh}
    \draw[gray!55] (\xs,0) -- (\xe,\ye);
  }
  \draw[very thick, ink] (0,0) -- (4,0) -- (2,{4*\hh}) -- cycle;
  \node[below left] at (0,0) {$B$};
  \node[below right] at (4,0) {$C$};
  \node[above] at (2,{4*\hh}) {$A$};
  \coordinate (d1) at (2,0);
  \coordinate (d2) at ({1+0.5},{\hh});
  \coordinate (d3) at ({2+0.5},{\hh});
  \fill[armRed!30] (d1) -- (d2) -- (d3) -- cycle;
  \draw[armRed!75!black, thick] (d1) -- (d2) -- (d3) -- cycle;
  \node[armRed!75!black] at (3.55,1.0) {շրջված};
  \draw[armBlue!75!black, very thick]
     (1,0) -- (3,0) -- (2,{2*\hh}) -- cycle;
  \node[armBlue!75!black] at (2,-0.35) {ուղիղ, $2$ կողմ};
\end{tikzpicture}
\end{center}

\textbf{Քանակներն ըստ չափի $n=8$-ի համար (ինչ որ իրականում գումարեցինք).}
\begin{center}
\small
\begin{tabular}{@{}c|cccccccc|c@{}}
\toprule
կողմ $k$ & 1 & 2 & 3 & 4 & 5 & 6 & 7 & 8 & ընդամենը\\
\midrule
ուղիղ    & 36 & 28 & 21 & 15 & 10 & 6 & 3 & 1 & \textbf{120}\\
շրջված   & 28 & 15 & 6  & 1  & 0 & 0 & 0 & 0 & \textbf{50}\\
\bottomrule
\end{tabular}
\end{center}

% ============================================================================
\section*{Խնդիր 12 - Գծային հավասարման սահմանափակ ամբողջ լուծումները}

\begin{problembox}[12]
Գտիր համակարգի ամբողջ լուծումների քանակը՝
$$x_1+x_2+x_3=40,\qquad 4\le x_1\le 15,\quad 9\le x_2\le 18,\quad 5\le x_3\le 16.$$
\end{problembox}

\begin{ideabox}
Երկու խոչընդոտ կա մեր և պարզ «աստղեր և ձողիկներ» հաշվի միջև՝ փոփոխականներն ունեն \textbf{ստորին}
սահմաններ (չեն սկսում $0$-ից) և \textbf{վերին} սահմաններ (չեն կարող անսահման աճել)։ Դրանք հեռացնում
ենք մեկ առ մեկ։

\textbf{Ստորին սահմանները ձրի են՝} տեղադրիր $y_i=x_i-(\text{իր մինիմումը})$, որպեսզի յուրաքանչյուր
$y_i\ge 0$ լինի և ընդհանուրը համապատասխանաբար նվազի։ \textbf{Վերին սահմանները պահանջում են
կցում-արտաքսում՝} նախ հաշվիր լուծումները՝ անտեսելով վերին սահմանները (հեշտ աստղեր և ձողիկներ), ապա
հանիր «անօրինական» լուծումները, որտեղ ինչ-որ $y_i$ գերազանցում է իր սահմանը։ Խելացի մասը՝ «$y_i$-ն
գերազանցում է իր $c_i$ սահմանը» նշանակում է «$y_i\ge c_i+1$», ինչը կիրառում ենք՝ նախապես $c_i+1$
միավոր հանձնելով $y_i$-ին և մնացածը ազատ հաշվելով։ Միանգամից երկու փոփոխականի գերազանցումները հանվում
են երկու անգամ, ուստի ԿԱՍ-ը դրանք հետ է ավելացնում։ Ամբողջ լուծումը մեկ կոկիկ նշանափոխ գումար է։
\end{ideabox}

\textbf{Քայլ 1 - Վերացնենք ստորին սահմանները։}
Դիցուք $y_1=x_1-4,\ y_2=x_2-9,\ y_3=x_3-5$։ Այդ դեպքում $y_i\ge 0$ և
$$y_1+y_2+y_3=40-4-9-5=22,\qquad 0\le y_1\le 11,\ \ 0\le y_2\le 9,\ \ 0\le y_3\le 11.$$
(Սահմանները՝ $15-4=11$, $18-9=9$, $16-5=11$։)

\textbf{Քայլ 2 - Հաշվենք՝ անտեսելով վերին սահմանները։}
$y_1+y_2+y_3=22$-ի ոչ բացասական լուծումներն աստղերով և ձողիկներով՝
$$\binom{22+3-1}{3-1}=\binom{24}{2}=276.$$

\textbf{Քայլ 3 - Հանենք միայնակ գերազանցումները։}
«$y_i$-ն իր սահմանից վեր» նշանակում է $y_i\ge \text{սահման}+1$՝ տուր նրան այդքան միավոր և հաշվիր
$22-(\text{սահման}+1)$ կրճատված ընդհանուրի լուծումները՝
$$|y_1\ge 12|=\binom{(22-12)+2}{2}=\binom{12}{2}=66,\quad
  |y_2\ge 10|=\binom{(22-10)+2}{2}=\binom{14}{2}=91,$$
$$|y_3\ge 12|=\binom{(22-12)+2}{2}=\binom{12}{2}=66.$$

\textbf{Քայլ 4 - Հետ ավելացնենք կրկնակի գերազանցումները։}
Միանգամից երկու սահման գերազանցված (տուր երկուսին էլ իրենց $\text{սահման}+1$ միավորները, ապա հաշվիր
մնացածը)՝
$$|y_1\ge12,\,y_2\ge10|=\binom{(22-12-10)+2}{2}=\binom{2}{2}=1,\quad
  |y_2\ge10,\,y_3\ge12|=\binom{2}{2}=1,$$
$$|y_1\ge12,\,y_3\ge12|:\ 22-12-12=-2<0\ \Rightarrow\ 0.$$
Միանգամից բոլոր երեք սահմանները պահանջում են $12+10+12=34>22$ միավոր, անհնար է, ուստի եռյակ անդամը
$0$ է։

\textbf{Քայլ 5 - Հավաքենք (կցում-արտաքսում)։}
$$276-(66+91+66)+(1+1+0)-0=276-223+2=55.$$

\textbf{Ստուգում։} $x_1\in[4,15],\,x_2\in[9,18],\,x_3\in[5,16]$ վրայով ուղիղ եռակի ցիկլը, որը հաշվում
է $x_1+x_2+x_3=40$ բավարարողները, վերադարձնում է ճիշտ $55$։ \checkmark

\answer{$\displaystyle \binom{24}{2}-\Big[\binom{12}{2}+\binom{14}{2}+\binom{12}{2}\Big]+\big[1+1\big]=55$}

\begin{methodbox}
\textbf{Փոխանցելի տեխնիկա - սահմանափակ կոմպոզիցիաներ՝} $x_1+\dots+x_k=N$-ի ամբողջ լուծումները
$a_i\le x_i\le b_i$ պայմանով հաշվելու համար՝
\begin{enumerate}[nosep,leftmargin=2em]
  \item Տեղադրիր $y_i=x_i-a_i$՝ բոլոր ստորին սահմանները $0$ դարձնելու համար (և $N$-ը նվազեցրու
        $\sum a_i$-ով)։
  \item Հաշվիր առանց վերին սահմանների՝ $\binom{N'+k-1}{k-1}$։
  \item ԿԱՍ-ով հանիր/ավելացրու անդամներ, որտեղ փոփոխականների ընտրված $S$ ենթաբազմությունը
        գերազանցում է՝ յուրաքանչյուր այդպիսի փոփոխական ընդհանուրից հանում է $(\text{սահման}_i+1)$,
        $(-1)^{|S|}$ նշանով։ Բաց թող ցանկացած անդամ, որի մնացած ընդհանուրը բացասական է դառնում (այն
        տալիս է $0$)։
\end{enumerate}
\end{methodbox}

\medskip
\textbf{ԿԱՍ-ի հաշվառումը մի հայացքով։}
\begin{center}
\small
\renewcommand{\arraystretch}{1.15}
\begin{tabular}{@{}llrcr@{}}
\toprule
գերազանցման բազմություն & մնացած ընդհանուր & $\binom{\cdot+2}{2}$ & նշան & ներդրում\\
\midrule
ոչ մեկը               & $22$ & $276$ & $+$ & $+276$\\
$y_1\ge12$           & $10$ & $66$  & $-$ & $-66$\\
$y_2\ge10$           & $12$ & $91$  & $-$ & $-91$\\
$y_3\ge12$           & $10$ & $66$  & $-$ & $-66$\\
$y_1\ge12,\,y_2\ge10$& $0$  & $1$   & $+$ & $+1$\\
$y_2\ge10,\,y_3\ge12$& $0$  & $1$   & $+$ & $+1$\\
$y_1\ge12,\,y_3\ge12$& $-2$ & $0$   & $+$ & $0$\\
բոլոր երեքը           & $-12$& $0$   & $-$ & $0$\\
\midrule
\multicolumn{4}{r}{\textbf{զուտ ընդամենը (ներդրումների գումար)}} & $\mathbf{55}$\\
\bottomrule
\end{tabular}
\end{center}

\end{document}
